\documentclass[10pt]{article}

\usepackage[utf8]{inputenc}

\usepackage[T1]{fontenc}

\usepackage{amsmath}

\usepackage{amsfonts}

\usepackage{amssymb}

\usepackage[version=4]{mhchem}

\usepackage{stmaryrd}

\usepackage{mathrsfs}

\usepackage{CJKutf8}

\usepackage{bbold}



\begin{document}

ориентированный на задачи преобразования символьной информации; в первоначальном варианте назывался «метаалгоритмическим языком» (см. [1]). Р. был создан как универсальный метаязык для описания преобразований языковых объектов. Он используется для трансляции с одного алгоритмич. языка на другой, для машинного выполнения аналитич. выкладок, доказательства теорем, перевода с естественных языков и т. П.



Запись алгоритма на $\mathrm{P}$. представляет описание нек-рого числа рекурсивных функций на множестве в ыр а ж е и й (т. е. последовательностей символов и скобок), правильно построенных (в обычном смысле) относительно скобок. Значение функции $\varphi$ при аргументе £ изображается на Р. в виде $\boldsymbol{К} \varphi$ зн а к к о н к р е т и а ц и и, служащий для явного указания на необходимость вычисления значения функции, а символ - означает закрывающую скобку для $\boldsymbol{K}$. Описание функции распадается на несколько предложений (п р а в и л к о н к р е т и 3 а ц и и), относящихся к случаям, когда аргумент имеет тот или иной частный вид. Напр., функция сложения в рекурсивной арифометике описывается на Р. двумя предложениями:



\begin{enumerate}

  \item $\boldsymbol{K}+(\boldsymbol{E} A)(0) \sqsubseteq \boldsymbol{E} A$,



  \item $\boldsymbol{K}+(\boldsymbol{E A})(\boldsymbol{E B 1}) \sqsupseteq \boldsymbol{K}+(\boldsymbol{E} A)(\boldsymbol{E} B) \perp 1$.



\end{enumerate}



Предложение состоит из левой и правой частей, разделяемых знаком д, и может включать свободные переменные (в примере әто - $\boldsymbol{E} A$ и $\boldsymbol{E} B$ ). Оно считается применимым для конкретизации выражения вида Кчб -, если это последнее может быть отождествлено с левой частью предложения при нек-рых значениях входящих в нее свободных переменных. Применение предложения состоит в замене конкретизируемого выражения на правую часть предложения, в к-рой свободные переменные замецены их значениями. Для вычисления значения функции предложения рассматриваются последовательно, и применяется первое из них, оказавшееся подходящим. Этот процесс повторяется, пока в объект работы входят знаки $\boldsymbol{K}$.



Для реализации программ на Р . разработаны әффективные трансляторы (см. [3], [4]; пример использования Р. для мапинного выполнения выкладок в теоретич. фиизике см. в [5]).



Лum.: [1] Т у р ч и н В. Ф., “Кибернетика», 1968, № 4, с. $45-54 ;[2]$ Т у ч ч и н В. Ф., С е р д о б о л ь с к и й В.И.,

«Кибернетика", 1969 , № 3 , с. $58-62 ;[3] \Phi$ л о р н ц е в С. Н., “Кибернетика", 1969 , № 3 , с. $58-62 ;[3] \Phi$ л о р е н ц е в С. Н.,

О л ю н и н В. Ю., Т у р ч и н В. Ф., «Тр. 1 Всесоюзн. конфеО л ю н и н В. Ю., Т у р ч и н В. Ф., “Тр. 1 Всесоюзн. конфе-

ренции по программированию", К., 1968 , с. $114-33 ;$ [4] P oренции по программированию", К., 1968, с. $114-33 ;$ [4] P о-

м а н е н к о С. А., Т у р ч и н В. Ф., «Тр. $2-\bar{n}$ Всесоюзн. конференции по программированию", Новосиб., 1970, с. $31-42$; [5] Б у д н и к А. П. [и д р.], "Ядерная физика", 1971, т. 14, c. $304-13$ B. Ф. Турчин.

В.



РЕФЛЕКСИВНОЕ ПРОСТРАНСТВО - банахово пространство $X$, совпадающее при каноническом вложении со своим вторым сопряженным $X^{* *}$. Подробнее, пусть $X^{*}$ - пространство, сопряженное с $X$, то есть совокупность всех непрерывных линейных функционалов, определенных на $X$. Если $\langle x, f\rangle$ - значение функционала $f \in X^{*}$ на әлементе $x \in X$, то при фиксированном $x$ и $f$, пробегающем $X^{*}$, выражение $\langle x, f\rangle=\mathscr{F}_{x}(f)$ будет линейным функционалом на $X^{*}$, то есть әлементом пространства $X^{* *}$. Пусть $\tilde{X} \subset X^{* *}-$ множество таких функционалов. Соответствие $x \rightarrow \mathscr{F}_{x}$ есть изоморфизм, не меняющий нормы $\|x\|=\left\|\mathcal{F}_{x}\right\|$. Если $\tilde{X}=X^{* *}$, то пространство $X$ наз. р е ф л е к с и в н ы м. Пространства $l_{p}$ и $L_{p}(a, b), p>1$, рефлексивны, пространство $C[a, b]$ не рефлексивно.



Пространство $X$ рефлексивно тогда и только тогда, когда $X^{*}$ рефлексивно. Другим критерием рефлексивности банахова пространства $X$ является слабая компактность единичного шара этого пространства.



Р. п. слабо полно, и замкнутое подпространство Р . п. рефлекси вно. Понятие рефлексивности естественным образом распространяется на локально выпуклые пространства.



Лит.: [1] Д а н ф о р д Н., Ш в а р ц Д ж., Линейные операторы, ч. 1 - Общая теория, пер. с англ., М., 1962; [2] И ос и Да К., Функциональный анализ, пер. с англ., М., 1967 โ3] К а н т о р о в и ч Л. В., А к и л о в Г. П., Функциональ-

ный анализ, 2 изд., М., 1977.



РЕФЛЕКСИВНОСТЬ - одно из свойств бинарных отношений. Отношение $R$ на множестве $A$ наз. р е фл е к с и в н ы м, если $a R a$ для любого $a \in A$. Примерами рефлексивных отношений являются равенство, эквивалентность, порядок. Т. С. Фодјанова.



РЕФЛЕКТИВНАЯ ПОДКАТЕГОРИЯ - подкатегория, содержащая «наибольшую» модель любого объекта категории. Точнее, полная подкатегория (5, категории Л наз. р е ф л е к т и в н о й, если (5) содержит (5-рефлектор (см. Рефлектор) для любого объекта категории. Полная подкатегория (5) категории $\mathscr{\Re}$ рефлективна тогда и только тогда, когда функтор вложения Id๕, $\mathfrak{\Re}:(5 \rightarrow \Re$ обладает сопряженным слева функтором $S: \mathscr{f} \rightarrow(5$. Функтор $S$ сопоставляет каждому объекту $A$ из $\mathfrak{\Re}$ его (5-рефлектор $S(A)$; морфизмы $\pi_{A}: A \rightarrow S(A)$, входящие в определение (5-рефлектора, определяют естественное преобразование тождественного функтора $\operatorname{Id} \Re$ в композицию функторов $S$ Idঞ, \begin{CJK}{UTF8}{mj}凡\end{CJK}. Двойственным к понятию Р. п. является понятие к о р еф ле к т и вн $о$ й п о д к а т е г о р и и.



Р. п. (5) наследует многие свойства объемлющей категории Я. Напр., морфизм $\mu \in$ Mor (5 тогда и только тогда является мономорфизмом в (5, когда он мономорфизм в Я. Поэтому всякая Р. п. локально малой слева категории локально мала слева. Р. п. обладает произведениями тех семейств объектов, для к-рых произведение существует в самой категории, при этом оба произведения оказываются изоморфными. То же самое справедливо и для любых пределов. С другой стороны, функтор $S$ переводит копределы из $\mathfrak{R}$ в копределы в $(5$. Поэтому Р. п. полной (слева) категории является полной (слева) категорией.



Пусть $\Re$ - полная локально малая категория. Всякая полная подкатегория (5) категории $\mathfrak{R}$, замкнутая относительно произведений и подобъектов своих объектов и содержащая правый нуль, является Р. п. В частности, всякое многообразие категории $\Re$ есть $\mathrm{P}$. п.



(5-рефлектор произвольного объекта $A$ строится следующим образом. Выбираются представители $\left(v_{i}, A_{i}\right)$, $i \in I$, таких факторобъектов объекта $A$, что $A_{i} \in \mathrm{Ob}(5$. Произведение



\[

P=\Pi_{i \in I} A_{i}\left(\pi_{i}\right)

\]



принадлежит (5, и (5-рефлектор $S(A)$ является образом однозначно определенного морфизма $\gamma: A \rightarrow P$, для к-рого $\gamma \pi_{i}=v_{i}, i \in I$.



П р и м е р ы. 1) Пусть $R$ - область целостности. Полная подкатегория инъективных модулей без кручения является Р. П. категории $R$-модулей без кручения: рефлекторами являются инъективные оболочки модулей. В частности, подкатегория полных абелевых групп без кручения есть Р. П. категории абелевых групп без кручения.



\begin{enumerate}

  \setcounter{enumi}{1}

  \item Полная подкатегория нормальных топологич. пространств есть Р. П. категории вполне регулярных топологич. пространств: рефлекторы строятся с помощью компактификации Чеха.



  \item Полная подкатегория пучков есть Р. п. категории предпучков: рефлекторы определяются функтором ассоциированного пучка. М. Ш. Цаленко.



\end{enumerate}



РЕФЛЕКТОР о б ъ е к т а к а т е г о р и и - понятие, описывающее «наибольшую» модель данного объекта в нек-ром классе объектов. А именно, пусть (5подкатегория категории $\mathfrak{\pi} ;$ объект $B \in(5$ наз. $\mathrm{p}$ е фле к т о ром о б в е к т а $A \in \mathbb{\Omega}$ в (5, или (5- р е ф-





\end{document}